\section{Results}\label{sec:results}

\subsection{The leakage-free fold: discrimination survives, the
federation gap does not}

Table~\ref{tab:leakfree} compares the five headline arms on the
leakage-free fold against the shipped in-partition numbers. Three
observations follow. First, \emph{discrimination generalises}: every
arm attains at least 0.906 ROC-AUC on the unseen partition, and the
task is easier than the in-partition protocol suggested (logistic
regression reaches 0.978/0.455). Second, the
centralised-versus-federated gap of the in-partition protocol
\emph{disappears}: the pooled MLP of identical architecture reaches
0.973/0.382 against FedProx's 0.976/0.308 --- a 0.003 ROC difference
and a reversed PR ordering. The in-partition gap (0.888 vs 0.954) was
an artefact of the protocol pairing, not a federation effect. Third,
the FedProx-versus-FedAvg contrast only partially survives: FedProx
leads by 0.070 ROC-AUC but trails FedAvg by 0.062 PR-AUC, and the
next subsection shows why neither difference should be read as an
algorithm property while BatchNorm sits in the loop.

\begin{table}[t]
\centering
\caption{Leakage-free fold (train partitions 1--4, single-pass test
on partition 5, natural 1.31\% prevalence) against the shipped
in-partition numbers. Threshold-conditional calibration metrics use
the validation-frozen prior-shift pipeline; ROC/PR are
threshold-free.}
\label{tab:leakfree}
\small
\begin{tabular}{@{}l c c c c c c c@{}}
\toprule
 & \multicolumn{2}{c}{\textbf{In-partition}} &
 \multicolumn{5}{c}{\textbf{Leakage-free fold}} \\
\cmidrule(r){2-3} \cmidrule(l){4-8}
\textbf{Model} & ROC & PR & ROC & PR & Brier & ECE & R@2\%FPR \\
\midrule
Logistic regression & 0.824 & 0.235 & 0.978 & 0.455 & 0.375 & 0.477 & 0.695 \\
XGBoost (centralised) & 0.977 & 0.493 & 0.971 & 0.457 & 0.055 & 0.101 & 0.692 \\
Centralised MLP & 0.888 & 0.153 & 0.973 & 0.382 & 0.544 & 0.687 & 0.555 \\
FedAvg MLP & 0.575 & 0.199 & 0.906 & 0.370 & 0.127 & 0.337 & 0.652 \\
FedProx MLP & 0.954 & 0.307 & 0.976 & 0.308 & 0.271 & 0.508 & 0.631 \\
\bottomrule
\end{tabular}
\end{table}

The calibration and threshold-transfer failures persist
leakage-free --- the balanced-validation carve (48.9\% prevalence)
again faces the natural-prevalence test (1.31\%), the same 26-fold
prior shift without the instance overlap --- so they are genuine
properties of the balanced-training/natural-test interface rather
than overlap artefacts: the validation-frozen FPR operating points
are violated by every arm (XGBoost realises 15.3\% FPR at the frozen
2\% target; FedProx 49.4\%), and only XGBoost's posteriors remain
usable (Brier 0.055 versus 0.127--0.544 for the neural arms).

\paragraph{Persistence baselines set the deployability floor.}
The toolkit of Leka et al.~\cite{leka2019} --- TSS, HSS,
persistence, climatology --- accompanies every operating point, and
the persistence baselines set two floors that are not decorative.
Same-region previous-window persistence (label inertia) achieves
TSS 0.967 / HSS 0.970 on the partition-5 test with no learned model
at all: at a 24-hour label horizon on $\sim$60-minute windows, an
active region that is flaring is still flaring, and \emph{no arm in
this paper --- XGBoost included --- beats that floor}. The
24-hour-lagged variant (forecast = the label of the same-region
window nearest $t-24$\,h) scores TSS 0.418 / HSS 0.434; on TSS,
XGBoost (0.848) and logistic regression (0.489) are the only arms
above it, and every neural arm at its frozen threshold sits below
it, while on HSS no arm clears the floor at all (XGBoost's HSS is
0.201). This is the classic Leka-era result, and a benchmark audit
is the right vehicle for stating it plainly: at this window geometry
the honest question is not whether models beat persistence, but
whether anything beats \emph{lagged} persistence --- and on this
fold, at TSS, two arms do, none at HSS. A committed sensitivity sweep
over eight candidate lag definitions agrees: the seven point-lag
variants land at TSS 0.392--0.418 with the same two arms clearing
each, while a lookback-flare variant scores 0.964 --- label inertia
under another name, which no arm clears.

\paragraph{Event-level evaluation.}
Regrouping the partition-5 test decisions into physical flare events
(65 true M/X-class events, validation-frozen thresholds, 1-hour
alert cooldown; Table~\ref{tab:eventlevel}), all five arms detect
effectively every event (64--65 of 65) with a median first-alert
lead of 23.4\,h before flare peak. That lead must be read for what it
is: the label is positive for every window in the 24 hours before a
flare, so a median first alert near the horizon's edge measures the
label boundary, not precursor physics. What separates the arms is
their false-alarm burden: XGBoost emits 1.7 false-alarm windows per
day versus 11.6--12.7 for the neural arms --- a $7\times$ difference
that mirrors the window-level calibration gap and confirms that the
deployability boundary on this benchmark is calibration, not
detection.

\begin{table}[t]
\centering
\caption{Event-level metrics on the leakage-free fold (65 true
M/X-class events, validation-frozen $F_2$ thresholds, 1-hour alert
cooldown). FA = false alarm.}
\label{tab:eventlevel}
\small
\begin{tabular}{@{}l c c c c@{}}
\toprule
\textbf{Model} & \textbf{Events detected} & \textbf{FA windows/day} &
\textbf{Alerts/day} & \textbf{Lead to peak (h)} \\
\midrule
Logistic regression & 65/65 (100\%) & 8.6 & 8.8 & 23.4 \\
XGBoost (centralised) & 64/65 (98.5\%) & 1.7 & 1.9 & 23.3 \\
Centralised MLP & 65/65 (100\%) & 12.7 & 12.8 & 23.4 \\
FedAvg MLP & 65/65 (100\%) & 11.6 & 11.8 & 23.4 \\
FedProx MLP & 65/65 (100\%) & 11.9 & 12.1 & 23.4 \\
\bottomrule
\end{tabular}
\end{table}

\subsection{The raw substrate: an apparent federated collapse}
\label{sec:rawcollapse}

Retraining the identical frozen protocol on the raw, unbalanced
benchmark separates the arms decisively (Table~\ref{tab:rawsubstr}).
The centralised arms are substrate-robust (logistic regression
unchanged to the third decimal; the centralised MLP's PR-AUC
\emph{rises} from 0.382 to 0.445; XGBoost's falls from 0.457 to
0.329 --- the synthetic positives were inflating its minority-class
precision, not teaching it). The federated MLP arms are not: FedAvg
falls to 0.875/0.056 with its calibrated posteriors never crossing
the validation-frozen threshold (test recall 0.000), FedProx to
0.930/0.149, and FedAvg detects \emph{zero} of the 65 events. Read
naively, this is a federation failure. The next two subsections show
it is mostly a BatchNorm artefact.

\begin{table}[t]
\centering
\caption{Raw-substrate retraining (train raw partitions 1--4,
natural 2.05\% prevalence, imputation and normalisation reproduced
in-pipeline with train-only parameters; single-pass test on raw
partition 5, natural 1.31\% prevalence; identical frozen protocol).
The leakage-free cleaned-substrate fold is reproduced for
comparison.}
\label{tab:rawsubstr}
\small
\begin{tabular}{@{}l c c c c c c@{}}
\toprule
 & \multicolumn{2}{c}{\textbf{Fold (cleaned)}} &
 \multicolumn{4}{c}{\textbf{Raw substrate}} \\
\cmidrule(r){2-3} \cmidrule(l){4-7}
\textbf{Model} & ROC & PR & ROC & PR & Brier & R@2\%FPR \\
\midrule
Logistic regression & 0.978 & 0.455 & 0.978 & 0.448 & 0.053 & 0.732 \\
XGBoost (centralised) & 0.971 & 0.457 & 0.974 & 0.329 & 0.018 & 0.587 \\
Centralised MLP & 0.973 & 0.382 & 0.971 & 0.445 & 0.010 & 0.596 \\
FedAvg MLP & 0.906 & 0.370 & 0.875 & 0.056 & 0.013 & 0.063 \\
FedProx MLP & 0.976 & 0.308 & 0.930 & 0.149 & 0.032 & 0.343 \\
\bottomrule
\end{tabular}
\end{table}

\subsection{BatchNorm finding 1: the in-partition collapse is a
composite evaluation artefact}

The client MLPs contain three \texttt{BatchNorm1d} layers, and local
training updates their running statistics every batch --- but the
federation transports \emph{parameters only}: running statistics are
buffers, not parameters, and buffers never move. The server-side
global model therefore evaluates with its initialisation buffers
forever (the shipped checkpoints carry \texttt{num\_batches\_tracked}
$=0$, \texttt{running\_mean} $=0$, \texttt{running\_var} $=1$ after
50 rounds), where each BN layer degenerates to the fixed affine map
$x\gamma+\beta$. The second half of the mechanism is the saturated
monitor: with init buffers, evaluated probabilities squeeze into a
narrow band, validation $F_1$ pins to its all-positive closed form
$2p/(1+p)$, and the best-validation checkpoint rule --- acting on
that saturated signal --- restored the round-20 FedAvg state over the
healthy round-5 state.

Table~\ref{tab:bndiag} decomposes the published in-partition FedAvg
numbers: the frozen 50-round protocol was re-run faithfully and the
identical global weights evaluated at three rounds under four BN
treatments. The verdict is unambiguous. The weights never lost
discriminative content: under evaluation-time batch statistics the
round-50 state --- whose shipped evaluation is 0.098, below chance
--- holds 0.953 ROC-AUC, and the round-20 state the rule shipped
holds 0.951. The in-partition collapse to 0.575 is a composite
evaluation artefact: untransported normalisation statistics ---
a mechanism documented since FedBN~\cite{li2021fedbn} and analysed
in subsequent BatchNorm-under-federation
studies~\cite{wang2023bn,guerraoui2024bn,bnscaffold2024} ---
compounded by a saturated monitor and the checkpoint rule it feeds.
A controlled
no-BN re-training in-partition confirms the training-side reading:
with the three BN layers removed and nothing else changed, neither
FedAvg nor FedProx diverges, the validation monitor stays
informative, and the 0.379 ROC-AUC gap between the two shipped
artefacts shrinks to 0.002--0.036 at every round.

\begin{table}[t]
\centering
\caption{Test ROC-AUC / PR-AUC of the identical in-partition FedAvg
global weights under four BatchNorm treatments (frozen protocol,
seed 42). Round 20 is the state the shipped checkpoint rule
returned. Arm D is what a buffer-transporting implementation would
have evaluated.}
\label{tab:bndiag}
\small
\begin{tabular}{@{}l c c c@{}}
\toprule
\textbf{BN treatment} & \textbf{Round 5} & \textbf{Round 20} & \textbf{Round 50} \\
\midrule
A: init buffers (shipped) & 0.950 / 0.264 & 0.575 / 0.199 & 0.098 / 0.010 \\
B: recalibrated (pooled) & 0.786 / 0.044 & 0.893 / 0.099 & 0.881 / 0.086 \\
C: batch statistics & 0.946 / 0.328 & 0.951 / 0.329 & 0.953 / 0.288 \\
D: transported round buffers & 0.885 / 0.157 & 0.931 / 0.194 & 0.935 / 0.209 \\
\bottomrule
\end{tabular}
\end{table}

\subsection{BatchNorm finding 2: removing BatchNorm restores the
federated MLP on the raw substrate}

The decisive control applies the same intervention to the raw
substrate: retrain both federated MLP arms with the three
\texttt{BatchNorm1d} layers replaced by the identity, changing
nothing else (same frozen protocol, shards, budget, aggregation,
seed). Table~\ref{tab:rawnobn} reports the outcome. FedAvg gains
$+0.088$ ROC-AUC and $+0.310$ PR-AUC ($6.5\times$); FedProx gains
$+0.045$ and $+0.196$. In gap terms: FedAvg's ROC gap to the
centralised MLP closes from 0.096 to 0.008 (91\%) and its PR gap
from 0.389 to 0.079 (80\%); FedProx's ROC gap closes \emph{past}
zero ($-0.004$ residual) and its PR gap from 0.296 to 0.100 (66\%).
Both no-BN arms sit level with their BatchNorm-free sequence
counterparts on ROC-AUC (FedAvg 0.963 vs the sequence arm's 0.958;
FedProx 0.975 vs 0.970), and no-BN FedProx \emph{matches} the
centralised MLP on ROC-AUC (0.975 vs 0.971): the $-0.004$ residual
sits inside the retrain nondeterminism the replication record
documents for exactly these arms (gate deltas $-0.002$ to $+0.023$
ROC-AUC), so the single-seed reading is parity within noise, not
federated precedence. The raw-substrate collapse of the BN-carrying
MLP arms is therefore dominated by BatchNorm under federation ---
not by the flattened-feature encoder, and not by federation per se.
Two boundaries keep the claim honest: (a) PR-AUC, not ROC-AUC, is
where work remains (no-BN FedProx's 0.345 sits below the centralised
MLP's 0.445 and the sequence arm's 0.404) --- what these arms
license is the attribution, not a new state-of-the-art row; (b) the
level claims now carry a second seed: re-executed at seed 43
(initialisation and shard draw reseeded, validation carve
unchanged), the shipped selections move by at most 0.002 ROC-AUC
(FedAvg 0.963$\to$0.961, FedProx 0.975$\to$0.976) and at most
0.024 PR-AUC without overturning an ordering, while the thresholded
validation-$F_1$ monitor is the seed-sensitive layer (best values
0.53$\to$0.66 and 0.09$\to$0.30) --- at 2\% prevalence the rank
metrics, not the monitor, carry the conclusions
(Section~\ref{sec:discussion}).

\begin{table}[t]
\centering
\caption{No-BatchNorm control on the raw substrate (single
intervention: three \texttt{BatchNorm1d} layers $\to$ identity,
nothing else changed; reference rows read from the released
artefacts).}
\label{tab:rawnobn}
\small
\begin{tabular}{@{}l c c@{}}
\toprule
\textbf{Arm (raw substrate)} & \textbf{ROC-AUC} & \textbf{PR-AUC} \\
\midrule
FedAvg MLP (BatchNorm, published) & 0.875 & 0.056 \\
FedAvg MLP (no BatchNorm, this control) & 0.963 & 0.366 \\
FedProx MLP (BatchNorm, published) & 0.930 & 0.149 \\
FedProx MLP (no BatchNorm, this control) & 0.975 & 0.345 \\
\midrule
Centralised MLP (reference) & 0.971 & 0.445 \\
FedAvg LSTM (reference, no BatchNorm) & 0.958 & 0.305 \\
FedProx LSTM (reference, no BatchNorm) & 0.970 & 0.404 \\
\bottomrule
\end{tabular}
\end{table}

\paragraph{The validation divergence, honestly.}
FedAvg's no-BN test ROC-AUC \emph{declines} from 0.974 (round 5) to
0.963 (round 50) while its validation ROC-AUC and $F_1$ still
improve (0.974 to 0.977; 0.345 to 0.527), so the shipped
best-validation rule selects the weakest test-ROC round of its own
monitored trajectory. A prevalence difference cannot explain this
divergence: ROC-AUC is a rank statistic, invariant to class
prevalence, so the validation carve's 2.05\% against the test fold's
1.31\% moves thresholded quantities, not the ranking. Two candidate
mechanisms remain, both testable with the released region-disjoint
split generator: the random validation carve shares active regions
with the training shards (the same leakage class the provenance
audit documents benchmark-wide), and partition 5 is the temporally
latest fold, so late-round specialisation to partitions 1--4 can
register as validation gain and test loss at once. The
within-fold region-disjoint validation re-run of both arms is
prepared in the released code and is the first remaining compute
item; FedProx's trajectory, for contrast, climbs monotonically from
0.931 to 0.975 --- the proximal term's stabilisation becoming
visible precisely once BatchNorm is removed.
